\section{Historical diagnostics and unsuccessful routes}
\label{app:history}
These results explain the sequence of research decisions; they are not
additional independent replications of the final candidate. Earlier exposed
evaluation questions remain development data. We retain each original
candidate decision and do not replace a failed nominee with a better-looking
control. The following tables reproduce aggregate point estimates only.

\paragraph{Prediction error did not establish editing value.}
The original G2 experiment used 200 training, 80 tuning and 120 held-out
GSM8K questions, five generation seeds, two anchor graphs and 12 edges.
Its cross-seed squared-risk estimates were 0.001416882 for the edgewise constant
baseline, 0.001541436 for the first-order predictor, and 0.002195856 for the second-order
predictor. Thus the second-order model failed its frozen prediction gate.
These are cross-seed risk estimates, not ordinary empirical MSE values.
Subsequent D2 controls showed that penalizing the intercept can induce
spurious positive edge credit even for a graph-independent constant response.
This is a property of the original objective, not evidence that the original
implementation violated its contract. Removing the intercept penalty reduced
risk, but the fixed D2 nominee achieved 0.000 percentage points of editing
gain on the already exposed evaluation block. It edited 238 of 240 anchors
and therefore did not acquire useful abstention merely by reducing risk.
The D3 graph-only second-order nominee likewise failed its separate
drift-enrichment gate (point estimate -0.001736111).

\paragraph{Mechanism evidence is distinct from system quality.}
The controlled-path experiments C3 and C4 each used a new 96-question
development cohort and two generation seeds. Their topology-sensitive
credit filters passed, yet their frozen graph nominees were worse than
their single-model and independent-aggregation controls (Table~\ref{tab:prior}).
C4 varied receiver identity and access to raw context; the table reports
the preselected full-context Qwen receiver, without selecting a different
factorial cell after outcomes. C3 used a Mistral terminal receiver, whereas
C4's nominated terminal receiver was Qwen. Cohorts and operators differ;
the two rows do not estimate the effect of the receiver change.

\begin{table}[t]
\centering\small
\begin{tabular}{lrrrr}
\hline
Stage & Graph EM & Single Q & Indep. & $D$ \\
\hline
C3 & 35.94 & 51.04 & 57.29 & 0.0625 \\
C4 & 50.52 & 58.85 & 67.71 & 0.1042 \\
\hline
\end{tabular}
\caption{Historical within-cohort comparisons. EM is a percentage;
$D$ is the frozen topology-sensitive credit contrast in its original
proportion units. The independent aggregator is Mistral in C3 and Qwen
in C4. Nonzero question support is 27/96 and 21/96, respectively.
Each row averages two generation seeds on the same question cohort;
seeds are not additional independent question samples.}
\label{tab:prior}
\end{table}

\paragraph{Feedback must beat inexpensive structural controls.}
Post-hoc replay on exposed C3 outputs found that affected-edge refresh
could match a zero-new-feedback protection rule without improving it.
This diagnostic motivated including protection and no-initialization
static actions in the later contract; it does not confirm that contract.
An execution audit also found disconnected source calls in the old operator.
Eliminating those calls is ordinary dead-code elimination and is shared
by every applicable method in the new operator. Accounting savings from
old traces are not measured runtime speedups of the new implementation.
